\section{The anchored Abel identity}
\label{sec:anchored-identity}

\begin{definition}[Translated logarithmic primitive]
\label{def:translated-primitive}
For \(r\ge0\), define
\[
 S_N=\sum_{n=1}^Na_n,\qquad
 L_{r,N}=\sum_{n=1}^N\frac{a_n}{r+n},
 \qquad L_{r,0}=0.
\]
\end{definition}

\begin{theorem}[Exact anchored Tauberian identity]
\label{thm:anchored-identity}
For every complex sequence, every real \(r\ge0\), and every
\(N\ge1\),
\begin{align}
 S_N
 &=(r+N)L_{r,N}-\sum_{m=0}^{N-1}L_{r,m},
 \label{eq:anchored-abel}\\
 \frac{S_N}{N}
 &=
 \left(1+\frac rN\right)L_{r,N}
 -\frac1N\sum_{m=0}^{N-1}L_{r,m}.
 \label{eq:anchored-defect-identity}
\end{align}
\end{theorem}

\begin{proof}
Since
\[
 a_n=(r+n)(L_{r,n}-L_{r,n-1}),
\]
summation and a shift of index give
\begin{align*}
 S_N
 &=
 \sum_{n=1}^N(r+n)L_{r,n}
 -\sum_{m=0}^{N-1}(r+m+1)L_{r,m}\\
 &=(r+N)L_{r,N}-\sum_{m=0}^{N-1}L_{r,m},
\end{align*}
because \(L_{r,0}=0\).  Division by \(N\) proves the second
identity.
\end{proof}

\begin{definition}[Anchored backward defect]
\label{def:anchored-defect}
For a finite sequence \(L=(L_m)_{0\le m\le N}\) with \(L_0=0\), put
\begin{equation}
 \cT_{r,N}(L)
 :=
 \left(1+\frac rN\right)L_N
 -\frac1N\sum_{m=0}^{N-1}L_m.
 \label{eq:anchored-defect}
\end{equation}
The usual signed backward Ces\`aro defect is
\[
 \cT_{0,N}(L)=L_N-\frac1N\sum_{m<N}L_m.
\]
\end{definition}

\begin{corollary}[Exact decomposition of the translation cost]
\label{cor:translation-cost}
\begin{equation}
 \boxed{\quad
 \cT_{r,N}(L)
 =\cT_{0,N}(L)+\frac rN L_N.
 \quad}
\label{eq:translation-cost}
\end{equation}
Consequently \(S_N/N\to0\) along any family of triples
\((r,N,a)\) if and only if \(\cT_{r,N}(L_r)\to0\).
\end{corollary}

\begin{proof}
Subtract the two definitions and use
\cref{eq:anchored-defect-identity}.
\end{proof}

\begin{remark}[The anchor is part of the main term]
\label{rem:anchor-main-term}
Equation \eqref{eq:translation-cost} is exact.  There is no reason for
\((r/N)L_N\) to be small merely because \(L_N=o(1)\).  The relevant
terminal scale is \(N/r\) when \(r\gg N\).
\end{remark}

\begin{proposition}[Absolute oscillation bound]
\label{prop:absolute-anchored}
Define
\[
 \cA_N(L)=\frac1N\sum_{m=0}^{N-1}|L_N-L_m|.
\]
Then
\begin{equation}
 \left|\frac{S_N}{N}\right|
 \le \cA_N(L_r)+\frac rN|L_{r,N}|.
\label{eq:absolute-anchored}
\end{equation}
\end{proposition}

\begin{proof}
Write the unshifted defect as
\[
 \cT_{0,N}(L)=\frac1N\sum_{m<N}(L_N-L_m)
\]
and apply the triangle inequality in
\eqref{eq:translation-cost}.
\end{proof}
